\section{Introduction}

\begin{frame}
  \frametitle{How to use this template} 

  \begin{block}{Preparation}
    \begin{enumerate}
    \item
      \emph{Fork} this repository
    \item
      Install \LaTeX{} and \emph{latexmk}
    \end{enumerate}
  \end{block}

  \begin{block}{Build pdf}
    Run \raise1pt\hbox{\fbox{\lstinline[language = sh]|make|}}
    \begin{itemize}
      \thusitem slide.pdf should be generated
    \end{itemize}  
  \end{block}
\end{frame}

\section{Some Templates}

\begin{frame}[fragile]
  \frametitle{Your first slide}
  \begin{columns}
    \begin{column}{0.55\textwidth}
      \lstinputlisting[
        language = tex,
        linewidth = \textwidth
      ]{./fig/sample-02.tex}
    \end{column}
    \begin{column}{0.45\textwidth}
      \begin{itemize}
      \item
        You can set the font size{\footnotesize (Now it's 10 pt)}
      \item
        Start comments with \alert{\%}
      \item 
        Each pages are called \emph{frames} in Beamer
      \end{itemize}
    \end{column}
  \end{columns}

\end{frame}

\begin{frame}[fragile]
  \frametitle{Simple List}
  \begin{columns}
    \begin{column}{0.51\textwidth}
      \lstinputlisting[
        language = tex,
        linewidth = \textwidth
      ]{./fig/sample-03.tex}
    \end{column}
    \begin{column}{0.49\textwidth}
      Use \emph{itemize} to list things
      
      \begin{itemize}
      \item This is the first item. 
      \item This is the second item.
      \item This is the third item.
      \end{itemize}
    \end{column}
  \end{columns}

  \begin{columns}
    \begin{column}{0.51\textwidth}
      \lstinputlisting[
        language = tex,
        linewidth = \textwidth
      ]{./fig/sample-04.tex}
    \end{column}
    \begin{column}{0.49\textwidth}
      Use \emph{enumerate} if the order matters
      
      \begin{enumerate}
      \item Do this first
      \item And then do that
      \end{enumerate}
    \end{column}
  \end{columns}

\end{frame}

\begin{frame}[fragile]
  \frametitle{Bored with bullet points?}
  \begin{columns}
    \begin{column}{0.54\textwidth}
      \lstinputlisting[
        language = tex,
        linewidth = \textwidth
      ]{./fig/sample-05.tex}
    \end{column}
    \begin{column}{0.46\textwidth}
      \begin{itemize}
        \ngitem
        Bullet points cannot show 
        the relation of the items
        \thusitem You may find it good to use some extra symbols
      \end{itemize}
    \end{column}
  \end{columns}

  \begin{mybox}{Custom \texttt{\textbackslash item} Macros}
    \begin{itemize}
      \okitem \verb|\okitem|: OK !
      \ngitem \verb|\ngitem|: NG
      \thusitem \verb|\thusitem|: Thus, \dots
      \butitem \verb|\butitem|: But, \dots
      \egitem \verb|\egitem|: For example, \dots
    \end{itemize}
  \end{mybox}
\end{frame}

\begin{frame}[fragile]
  \frametitle{Emphasize your key points!}
  \begin{enumerate}
  \item
    Make sure to \verb|\emph{Emphasize}| your \emph{key points}
    \footnote{\insertfootnotemark I let them to be bold}
  \item
    \verb|\Emph{Use this}| to empasize \Emph{very important point}
    \footnote{\insertfootnotemark I defined customly}
  \item 
    We have \verb|\EMPH{this command}| for the \EMPH{very very very important notes}%
    \footnote{\insertfootnotemark I defined customly}
  \end{enumerate}
  
\end{frame}


\begin{frame}[fragile]
  \frametitle{Block, Definition, Examples}
  \begin{columns}
    \begin{column}{0.52\textwidth}
      \lstinputlisting[
        language = tex
      ]{./fig/sample-06.tex}
    \end{column}
    \begin{column}{0.48\textwidth}
      \begin{block}{Title of the block}
        Here is the body sentences of the block.
        Notice foo is not bar!
      \end{block}
    \end{column}
  \end{columns}
  \begin{columns}
    \begin{column}{0.52\textwidth}
      \lstinputlisting[
        language = tex
      ]{./fig/sample-07.tex}
    \end{column}
    \begin{column}{0.48\textwidth}
      \begin{definition}
        A prime number is \dots 
      \end{definition}

      \begin{example}   
        \begin{itemize}
          \okitem 2 is a prime number．
          \okitem 3 is also a prime number．
          \ngitem 4 is not a prime number．
        \end{itemize}
      \end{example}
    \end{column}
  \end{columns}
\end{frame}

\begin{frame}[fragile]{Source codes}
  \lstinputlisting[
    language = tex,
    linewidth = 0.7\textwidth
  ]{./fig/sample-10.tex}
  \begin{itemize}
  \item
    You need to pass \Emph{fragile} as the option of frame
  \item
    There are a lot of options in the listings package
    \begin{itemize}
      \egitem language, linewidth, $\dots$
    \end{itemize}
  \end{itemize}

\end{frame}

\begin{frame}[fragile]
  \frametitle{Inserting images}
  \begin{columns}
    \begin{column}{0.67\textwidth}
      \lstinputlisting[
        language = tex,
        linewidth = \textwidth
      ]{./fig/sample-11.tex}
    \end{column}
    \begin{column}{0.33\textwidth}
      \includegraphics[width = 3cm]{./fig/logo}
    \end{column}
  \end{columns}
\end{frame}

\begin{frame}{Equation}
  \begin{equation}
    x = a_0 + \cfrac{1}{a_1 
      + \cfrac{1}{a_2 
        + \cfrac{1}{a_3 + \cfrac{1}{a_4} } } }
  \end{equation}

  \[
  \sqrt[n]{1+x+x^2+x^3+\dots+x^n}
  \]
\end{frame}

